\section{Conclusion}\label{sec:conclusion}

The Bailout Protocol specifies the runtime mechanism by which the BX3 Framework's upstream accountability guarantee is enforced in real-time multi-agent execution. The contribution is architectural and formal: the protocol defines a deterministic state machine, triggered by the Bounds Engine's correct identification of the boundary of its own authority, that propagates every unresolvable exception state directly to a named human accountability anchor, bypassing all intermediate machine actors, with complete diagnostic context, and with a tamper-evident forensic record of the entire propagation chain.

The protocol's correctness properties were established formally. The Safety property (Theorem~1) guarantees that no machine actor may issue a terminal resolution to any BailoutTrigger: the state machine has no transition from any state other than \texttt{HUMAN\_ACKNOWLEDGED} to \texttt{RESOLVED}, and the Fact Layer hard block $\Gamma_B$ is monotonic, applied at \texttt{BOUNDED} and not released until \texttt{RESOLVED}. The Liveness property (Theorem~2) guarantees that every valid trigger event reaches a human anchor within a bounded number of escalation steps, because the system hierarchy is a finite rooted tree and the escalation algorithm ascends the parent chain strictly. The Accountability property (Theorem~3) guarantees that the human anchor identified at resolution time is the same human anchor designated at the node's provisioning time, established by the immutable parent chain and enforced by the Fact Layer's deterministic state machine.

These three properties jointly constitute the upstream accountability guarantee as a runtime invariant, not merely a design aspiration. The guarantee is not maintained by the reliability of any individual component, by operational policy, or by the vigilance of human monitors. It is maintained by the layer architecture of the BX3 Framework, enforced at the Fact Layer, and triggered by the Bounds Engine's boundary conditions of authority. The system fails upward into human consciousness---it never fails downward into algorithmic chaos.

The Bailout Protocol is the runtime expression of this guarantee. The BX3 Framework's three-layer model---Purpose Layer for human accountability, Bounds Engine for constrained execution, Fact Layer for deterministic enforcement---is the architectural substrate that makes the protocol structurally necessary, not merely recommended. Every mechanism in the protocol maps onto exactly one layer: the trigger condition detection maps onto the Bounds Engine's boundary evaluation function; the hard block maps onto the Fact Layer's deterministic enforcement; the escalation terminus maps onto the Purpose Layer's exclusive resolution authority. This mapping is not a design pattern. It is a structural necessity. Removing any layer causes the corresponding guarantee to collapse.

The Bailout Protocol was specified in sufficient formal detail to serve as an authoritative reference for implementation, audit, and formal verification. The state machine, trigger conditions, escalation algorithm, bypass mechanism, and timeout handling were each defined precisely. Limitations were acknowledged transparently: human response latency in real-time contexts, adversarial vulnerability of the escalation path, scalability under high node density, and exploitation risk from deliberate trigger manipulation. Three productive directions for future work were identified: mechanized formal verification, adaptive timeout policy, and distributed human anchor pools.

The broader contribution is to the problem of accountability in multi-agent AI systems. As autonomous agents are deployed in higher-stakes environments---agricultural operations, infrastructure management, healthcare, financial systems---the question of who is responsible when something goes wrong is no longer a philosophical or regulatory abstraction. It is an engineering requirement. The Bailout Protocol demonstrates that mandatory human accountability as an architectural fallback is not merely desirable but formally achievable: that a deterministic, auditable, structurally enforced protocol can guarantee that every unresolvable exception state in a multi-agent system terminates in human judgment, with complete forensic context, without probabilistic components, and without relying on human vigilance as the primary safety mechanism. This is the contribution that the BX3 Framework was designed to deliver.
